\section{Discussion}
\label{sec:discussion}

\subsection{Generality of Sign Misattribution}
\label{sec:discussion_generality}

The sign misattribution mechanism identified in Section~\ref{sec:m2_sign_misattribution} is not specific to RAGEN, ALFWorld, or GRPO. The failure arises from three structural properties that co-occur in any multi-turn agent RL setup: (a) step-level credit assignment via windowing, (b) sparse binary episode reward, and (c) group-relative advantage normalization. When these three conditions hold, the advantage sign for each window collapses to a proxy for episode outcome regardless of local action quality (Equation~\ref{eq:sign_correlation}). The environment details, the specific RL algorithm, and the model architecture are irrelevant to this structural argument.

Consider the broader landscape of LLM agent training. Web agents operating in WebArena~\cite{zhou2024webarena} execute 20 to 40 interaction steps per episode with binary task completion rewards. Scientific experimentation agents in ScienceWorld~\cite{wang2022scienceworld} follow multi-step experimental protocols where success is determined only at the final measurement. GUI agents in mobile and desktop environments issue sequences of clicks and keystrokes with success defined by whether the target UI state is reached. Tool-use agents chain API calls where only the final output is evaluated. All of these settings combine long trajectories with sparse terminal rewards, and any framework that decomposes these trajectories into independently optimized windows will exhibit sign misattribution.

Proposition~2 (Section~\ref{sec:formal_probability}) reveals a further implication: the severity of misattribution scales with task difficulty. As the baseline success rate $p$ decreases, $P(\text{correct} \mid \text{update}) \to p$ in the large-$K$ limit. Environments with longer horizons or more complex reasoning requirements tend to have lower baseline success rates, so sign misattribution becomes \textit{more} severe precisely where credit assignment is most needed. A web agent with $p = 5\%$ baseline success would experience correct gradient direction in only 5\% of updates, a regime where training cannot converge regardless of other hyperparameter choices.

\subsection{Direction Correctness vs.\ Credit Granularity}
\label{sec:discussion_tradeoff}

Episode-level GRPO eliminates sign misattribution by treating the full trajectory as a single optimization sample. The cost of this choice is that all tokens in the episode receive the same advantage. A correct action at step 3 and a redundant action at step 40 are reinforced with identical magnitude and sign. This represents a trade-off between two desiderata: direction correctness (whether the gradient sign is aligned with local action quality) and credit granularity (whether the gradient magnitude reflects per-step contribution to success).

At the low success rates typical of early training ($p = 12.5\%$ in our experiments), direction correctness dominates. When 80.9\% of gradient updates point in the wrong direction, the optimization actively destroys good behaviors. Uniform reinforcement of all steps in a successful episode is imprecise but at least preserves correct sub-strategies within the trajectory. The asymmetry is stark: coarse-but-correct gradients produce slow convergence, while fine-but-wrong gradients produce divergence.

As the policy improves and $p$ increases, the marginal value of per-step credit assignment grows. At $p = 50\%$, Proposition~2 predicts $P(\text{correct} \mid \text{update}) \approx 50\%$, meaning step-independent training would produce a roughly unbiased gradient in expectation (though with high variance). In this regime, the coarseness of episode-level training may become a binding constraint on sample efficiency. Methods that combine episode-level direction correctness with finer-grained credit represent a natural extension. GAE with episode-level correction (using a value function for magnitude estimation while constraining signs to be consistent with episode outcome) and GiGPO's nested group structure~\cite{feng2025gigpo} are two concrete approaches. The connection to GAE(1,1) established in Section~\ref{sec:classical_connection} suggests a principled interpolation: begin with $\lambda = 1$ (episode-level) during early training when direction correctness matters most, and anneal toward $\lambda < 1$ as the policy stabilizes and per-step credit becomes valuable.

\subsection{Relationship to Classical Credit Assignment}
\label{sec:discussion_classical}

Sign misattribution is a specific manifestation of the temporal credit assignment problem, adapted to the LLM agent setting. In classical RL, temporal credit assignment is addressed through value functions (TD learning), eligibility traces (TD($\lambda$)), or return decomposition methods. The solutions all share a common structure: they provide mechanisms for information to propagate across time steps within an episode, so that the contribution of each step to eventual reward can be estimated.

Step-independent windowing violates this structure in a way that is unusual even within the classical RL literature. TD(0) assigns credit through one-step bootstrapped returns, which is local but connected: the value of state $s_t$ depends on $s_{t+1}$, creating a chain that propagates information across the episode over multiple training iterations. Monte Carlo methods assign credit through full-episode returns, providing direct (if noisy) temporal coupling. Step-independent training creates \textit{disconnected} windows that share no temporal coupling, neither through bootstrapping nor through shared returns. Each window is optimized as if it were an independent bandit arm that happens to receive a reward correlated with the episode outcome. This falls below the TD(0)-to-MC spectrum entirely.

The multi-agent analogy from Section~\ref{sec:classical_connection} clarifies the structural novelty. In cooperative multi-agent RL, the credit assignment challenge of shared team reward is well-recognized, and solutions such as counterfactual baselines (COMA~\cite{foerster2018coma}) or value decomposition (QMIX~\cite{rashid2018qmix}) exist. In multi-turn LLM agent training, the same structural problem arises from an implementation choice rather than from the problem formulation itself. The multi-agent credit assignment problem is inherent to team games; the sign misattribution problem in agent RL is artificial, introduced by the windowing design and removable by episode-level training.

\subsection{Dual Mechanism Interaction}
\label{sec:discussion_interaction}

Information loss (Mechanism 1) and sign misattribution (Mechanism 2) are not independent. They interact through the success rate $p$, creating a feedback loop that accelerates collapse at small window sizes. Mechanism 1 restricts the representable strategy space, imposing a ceiling on achievable task success for a given $h$. This ceiling determines the maximum $p$ that the policy can sustain. Mechanism 2 then drives performance below this ceiling at a rate that depends on $p$ itself: lower $p$ produces more severe misattribution (Proposition~2), which further reduces $p$, which further increases misattribution severity.

This interaction creates a vicious cycle. At $h = 2$, the information-loss ceiling is low, so the policy cannot achieve high $p$ even with perfect optimization. The low $p$ triggers severe sign misattribution ($\approx$80.9\% wrong-direction gradients at $p = 12.5\%$), which erodes whatever competence the model retains within its truncated context. The erosion reduces $p$ further, intensifying the misattribution. The two-phase collapse observed empirically (strategic failure followed by formatting failure) reflects this cascade: Mechanism 2 first destroys task-level strategies (reducing $p$ to near zero), then continues to corrupt lower-level action representations that were initially robust.

At larger $h$ values, the ceiling is higher, sustaining a higher $p$ that partially mitigates misattribution severity. This predicts a more gradual degradation rather than catastrophic collapse. The h-sweep ablation (Section~\ref{sec:results_hsweep}) tests this prediction by examining whether larger $h$ values exhibit slower collapse trajectories. Episode-level training breaks the cycle entirely by eliminating both mechanisms simultaneously: the full context removes the information-loss ceiling, and the single episode-level advantage removes the opportunity for sign misattribution.


\subsection{Practical Implications}
\label{sec:practical_implications}

Our diagnosis suggests a simple decision procedure for practitioners training multi-turn agents with RL. The core question is whether the training framework preserves correct gradient direction across the full episode.

\begin{enumerate}
    \item \textbf{If using step-independent windowing} (standard PPO/GRPO with sliding windows, $h = 2$--$5$): the credit agreement rate collapses to chance level (52.07\%, Table~\ref{tab:sign_verification}), and the $\gamma$-sweep confirms that all discount factors produce collapse, including $\gamma = 1.0$ (Section~\ref{sec:sign_verification}). At 8B scale, restricting to $h = 2$ incurs a 37.5 percentage point performance gap from information loss alone (Section~\ref{sec:cross_scale}). These failures are structural, not tunable. \textit{Avoid step-independent windowing for tasks with sequential dependencies spanning more than $h$ steps.}

    \item \textbf{If using episode-level credit} (full trajectory as a single optimization sample): this eliminates both sign misattribution and information loss. At 1.7B, episode-level GRPO recovers from 0\% to 23.88\% OOD success; at 8B, it reaches 90.3\% OOD from a collapsing baseline (Table~\ref{tab:main_results}, Table~\ref{tab:8b_results}). The cost is uniform advantage across all tokens in the episode: correct and redundant actions receive identical reinforcement. For tasks where baseline success is low ($p < 0.3$), this coarseness is a minor concern because direction correctness dominates (Section~\ref{sec:discussion_tradeoff}). \textit{Use episode-level credit as the default for multi-turn agent RL.}

    \item \textbf{If episode-level training plateaus}: once the policy reaches moderate success rates ($p > 0.5$), the uniform advantage becomes the binding constraint on sample efficiency. At this point, finer-grained methods (turn-level credit, value-function-based decomposition, $\lambda$-annealing from $\lambda = 1$ toward $\lambda < 1$) can improve credit granularity without reintroducing sign errors. The key requirement is that any decomposition must preserve episode-level sign consistency. \textit{Adopt finer credit assignment only after the policy is stable enough that direction errors are rare.}
\end{enumerate}

The ordering matters: starting with fine-grained credit on a low-success-rate policy ($p < 0.15$) produces catastrophic collapse (Section~\ref{sec:results_main}), while starting with episode-level credit and refining later sacrifices only convergence speed. The asymmetry between wrong-direction and imprecise-magnitude gradients makes the conservative choice strictly safer.

\subsection{Limitations}
\label{sec:limitations}

\paragraph{Single environment.} Our experiments are conducted exclusively on ALFWorld, a text-based household environment. The dual mechanism analysis depends only on the training algorithm structure (windowing, binary reward, group normalization) and not on environment-specific properties, so the theoretical argument is environment-agnostic. However, empirical validation on additional multi-turn environments such as ScienceWorld~\cite{wang2022scienceworld} or WebShop~\cite{yao2022webshop} would strengthen the generality claim. Different environments may exhibit different baseline success rates, trajectory lengths, and dependency structures that modulate the practical severity of sign misattribution.

\paragraph{Single model scale.} All experiments use Qwen3-1.7B. The severity of sign misattribution depends on $p$ (Proposition~2), which varies with model capacity. A 7B or 14B model may achieve a higher SFT baseline, reducing the fraction of wrong-sign updates. The misattribution mechanism does not disappear at higher $p$ (it persists as long as $p < 1$), but its practical impact may be less catastrophic if the model starts closer to the ceiling. Scaling experiments are needed to characterize the relationship between model capacity, baseline success rate, and collapse dynamics.

\paragraph{Coarse credit in episode-level training.} Episode-level GRPO assigns uniform advantage to all tokens in an episode. For very long episodes ($T > 100$) or tasks with many independent sub-goals, this coarseness may limit learning efficiency even though it preserves direction correctness. The trade-off between direction correctness and credit granularity (Section~\ref{sec:discussion_tradeoff}) becomes increasingly relevant at higher success rates where the sign error is less severe but the efficiency cost of uniform reinforcement grows.

\paragraph{Selective verification not implemented.} The selective verification component (Section~\ref{sec:selective_verification}) remains a design proposal. We present the training methodology diagnosis and episode-level solution, but the downstream application of training a selective verification policy has not been implemented or evaluated. The contribution of this paper is the credit assignment analysis and fix, not the verification application.

\paragraph{Limited gradient analysis scope.} The empirical sign misattribution verification (Section~\ref{sec:sign_verification}) uses $h = 2$, representing the most extreme case of windowed training. Analysis across multiple $h$ values ($h \in \{5, 10, 20\}$) would reveal how misattribution severity varies with window size in practice and whether the theoretical predictions (which are $h$-independent for a given $p$) hold when different window sizes produce different effective $p$ values during training.
